\documentclass[letterpaper,12pt]{exam}
\usepackage{csc274videonotes}

\newcommand{\unit}{Unit 00}
\pagestyle{headandfoot}
\firstpageheader{CSC 274 \semester\ \  \unit}{}{Name: $\rule{6cm}{0.15mm}$}
\runningheader{CSC 274 \semester}{\unit}{Page \thepage\ of \numpages}
\firstpagefooter{}{}{}
\runningfooter{}{}{}

\begin{document}

\section*{Instructions}
\begin{itemize}
	\item Print out these pages.  If you do not have access to a printer, please contact the instructor.
	\item Handwrite your answers on the pages you print out. 
	   \begin{itemize}
		  \item \textbf{I will not accept just the answers written on a sheet of paper.}  
		  \item \textbf{I will not accept typed output unless it includes the questions.}
	   \end{itemize}
	\item If you have a flatbed scanner, please scan the images into a .pdf file.  
	  \begin{itemize}
         \item PDF files are much better than individual jpg files.
         \item If you don't have a flatbed scanner, please use a "camscanner" app on your phone to produce a single .pdf file for the assignment.
      \end{itemize}
\end{itemize}

\section*{Unit 00\_05 --- Videos and Notes} % Unnumbered section

\begin{center}
\textbf{This section of the notes goes with the video and notes Unit 00.05}

The notes I am lecturing from are at \url{https://github.com/noynaert/csc274repo/Unit_00_CourseOverview/Unit_00_05_Coursework.md}.
I probably will not put these links in the notes in the future.  You should be able to find lessons for other notes at 
\url{https://github.com/noynaert/csc274repo/} and then looking in the individual units.  Generally 1 unit is 1 week.
\end{center}

  \begin{itemize}
	\item The video notes are intended as \textbf{learning aids} and not grading exercises.
	\item Watch the videos and fill these pages out as you go.
	\item I will only spot check a few questions and look for completeness.  There are only a few ways you can lose points:
	   \begin{itemize}
	     \item Skip a question. 
		 \item Just copy some answer you found on the internet.  Use your own words.
	     \item Give a frivolous answer.  It is oK to say "I don't know" but if you do this you must also show that you have studied the question.  Explain what confuses you.  You may need to look up another source on the internet or on the handouts pages.
	     \item Cheat.  Don't use another student's answers, although it is fine (and encouraged) if you talk over specific questions with another student.  Just be sure to give your own answer and don't copy another student.  Also, don't use notes from another semester; I can be sneaky.
	   \end{itemize}
	\end{itemize}

If you don't know an answer try to figure it out.  Instead of an answer,
write what you do understand or why you are confused.  Maybe put a big \textcolor{BrickRed}{\Huge*} or \textcolor{BrickRed}{\Huge?} in the margin next to 
the question.

%\begin{questions} % Begins the questions environment
   
%\end{questions}
\section*{Unit 00\_10 --- Course Overview}

\begin{center}
At this point you should go to Canvas and watch the next video.

You may also go to the repository in github and find the handout sheet for the next video.  
The repo is at \url{https://github.com/noynaert/csc274repo/}
\end{center}
\begin{questions}
\question How are you filling out these notes?
\begin{checkboxes}
\choice I am writing the answers on a sheet of notebook paper which I will turn in (This is a wrong answer.)
\choice I am typing the answers (This is also a wrong answer)
\choice I am hand-writing the notes on printed copies of the .pdf file from the handouts. (This is the correct answer.)
\end{checkboxes}
\end{questions}

\includegraphics[width=1in]{tux}\label{end}

\end{document}